\subsection{G 的闭子群的空间}

我们用 \(\Sigma\) 表示 G 的闭子群组成的集合；若把每一测度 \(\alpha\in\Gamma\) 对应到子群 \(H_{\alpha}\)（即 \(\alpha\) 的支集），就得到 \(\Gamma\) 到 \(\Sigma\) 中的一个映射（称为典范映射），它显然是满射，并使我们可以典范地把 \(\Sigma\) 等同于 \(\Gamma\) 中比率为 \textgreater0 的位似群的同轨集。集合 \(\Sigma\)（赋有 \(\Gamma\) 上淡拓扑的商拓扑）称为 G 的闭子群的空间。

定理 1. --- 设 G 为局部紧群。G 的闭子群组成的空间 \(\Sigma\) 是紧的。此外，我们有下列性质：

\begin{enumerate}
\def\labelenumi{(\roman{enumi})}
\item
  \(\Sigma^0\)（\(\mathbf{G}\) 的单模闭子群组成的集合）在 \(\Sigma\) 中是闭的（从而是紧的）。
\item
  若 G 由 e 的一个紧邻域生成，则集合 \(\Sigma_{c}^{0}\)——由 G 中使得商空间 G/H 为紧的单模闭子群 H 组成——在 \(\Sigma^{0}\) 中是开的（从而是局部紧的）。
\item
  对 e 在 G 中的每一相对紧开邻域 U，集合 \(D_{U}\)——由 G 中使得 \(H \cap U = \{e\}\) 的离散子群 H 组成——在 \(\Sigma^{0}\) 中是闭的（从而是紧的）。
\end{enumerate}

由 No.~1 的命题 3 可知 \(\Sigma\) 同胚于 \(\Gamma_{\varphi}\)，从而由 No.~1 的命题 2 是紧的。此外，在 No.~2 开头已指出，集合 \(\Gamma^0\)——由 \(\alpha \in \Gamma\) 中使得 \(\mathbf{H}_{\alpha}\) 单模的测度组成——在 \(\Gamma\) 中是闭的；由于 \(\Gamma^0\) 在比率为 \(>0\) 的位似下稳定，像 \(\Sigma^0\)（\(\Gamma^0\) 在 \(\Sigma\) 中的）是 \(\Sigma\) 的一个闭子集，这就证明了 (i)。

性质 (ii) 将是下列命题的一个推论：

命题 6. --- 假设局部紧群 G 由 e 的一个紧邻域生成。则集合 \(\Gamma_{c}^{0}\)——由测度 \(\alpha\in\Gamma^{0}\)（使得 \(\mathbf{G} / \mathbf{H}_{\alpha}\) 为紧的）组成——在 \(\Gamma^0\) 中是开的，且 \(\Gamma_c^0\) 上映射 \(\alpha \mapsto \| \mu_{\alpha}\|\) 的限制是淡连续的。

用 No.~2 命题 5 的记号，我们有，对 \(g \in \mathcal{K}_{+}(\mathbf{G})\)，

\[
\Gamma^ {0} (g) \subset \Gamma_ {c} ^ {0}.\tag{6}
\]

因为，若 K 为 g 的支集，则关系 \(\int g(xs) d\alpha(s) \geqslant 1\)（对所有 \(x \in G\)）蕴含 \(KH_{\alpha} = G\)，因为积分在 \(KH_{\alpha}\) 的余集上显然为零，于是 \(G/H_{\alpha} = \pi_{\alpha}(K)\) 是紧的。给定一个测度 \(\alpha \in \Gamma_{c}^{0}\)，因此只需定义一个函数 \(g \in \mathcal{K}_{+}(G)\)，使得 \(\Gamma^{0}(g)\) 是 \(\alpha\) 在 \(\Gamma^{0}\) 中的一个邻域。由于 \(G/H_{\alpha}\) 是紧的，且典范映射 \(f \mapsto f_{\alpha}\)（从 \(\mathcal{K}_{+}(G)\) 到 \(\mathcal{K}_{+}(G/H_{\alpha})\) 中）是满射（Ch. VII, §2, No.~2），存在一个函数 \(g \in \mathcal{K}_{+}(G)\)，使得 \(\int g(xs) d\alpha(s) = 2\) 对每一 \(x \in G\) 成立。设 K 为 g 的（紧）支集，L 为 e 在 G 中的一个生成 G 的对称紧邻域；映射 \(\beta \mapsto g*\beta\)（从 \(\mathcal{M}_{+}(G)\) 到 \(\mathcal{C}(G)\) 中）是淡连续的（§4, No.~2 备注 1），故存在 \(\alpha\) 在 \(\Gamma^{0}\) 中的一个邻域 W，使得

\[
(g * \beta) (x) = \int g (x s) d \beta (s) \geqslant 1\tag{7}
\]

对所有 \(\beta \in W\) 与 \(x \in LK\) 成立。(7) 的左端在 \(\mathrm{KH}_{\beta}\) 之外为零，故关系 \(\beta \in W\) 蕴含

\[
\mathbf {L K} \subset \mathbf {K H} _ {\beta},
\]

由此通过对 n 作归纳推出，\(L^{n}K \subset KH_{\beta}\) 对每个整数 n \textgreater{} 0 成立；由于 L 生成 G，我们有 \(G = KH_{\beta}\)（对每一测度 \(\beta \in W\)），这就证明了 \(W \subset \Gamma_{c}^{0}\)。另一方面，(7) 的左端在 \(H_{\beta}\) 下右不变，故不等式 (7) 对 \(x \in LKH_{\beta} = G\) 也成立；因此 \(W \subset \Gamma^{0}(g)\)，这就证明了命题。

最后，(iii) 将是下列命题的一个推论：

命题 7. --- 设 \(N \subset \Gamma^{0}\) 为 G 的离散子群上的规范化 Haar 测度组成的子空间，且对 e 在 G 中的每一相对紧开邻域 U，设 \(N_{U}\) 为 N 中由 \(\alpha\)（满足 \(H_{\alpha} \cap U = \{e\}\)）组成的子集。则：

\begin{enumerate}
\def\labelenumi{\alph{enumi})}
\item
  \(\mathbf{N}_{\mathrm{U}}\) 是紧的。
\item
  \(\mathbf{N}_{\mathbf{U}}\) 在 \(\mathbf{N}\) 中的内部构成 \(\mathbf{N}\) 的一个覆盖（当 \(\mathbf{U}\) 取遍 \(e\) 在 \(\mathbf{G}\) 中的相对紧开邻域组成的集合时）。
\item
  要使 N 的一个子集 M 在 N 中相对紧，必须且只需存在 e 在 G 中的一个相对紧开邻域 U，使得 \(M \subset N_{U}\)。
\end{enumerate}

由于 \(D_{U}\) 是 \(N_{U}\) 在典范连续映射 \(\Gamma \rightarrow \Sigma\) 下的像，定理 1 的断言 (iii) 立即由命题 7 a) 得出。

为证明命题 7，我们注意到 \(N_{U}\) 可定义为 \(\Gamma^{0}\) 中由满足下列二条件的 \(\alpha\) 组成的子集

\[
\alpha (\{e \}) \geqslant 1 \quad \text { and } \quad \alpha (U) \leqslant 1.
\]

现在，若 A 在 G 中是紧的（相应地，是开的且相对紧的），则映射 \(\alpha \mapsto \alpha(A)\)（从 \(\mathcal{M}_{+}(G)\) 到 R 中）对淡拓扑是上半（相应地，下半）连续的（Ch. IV, §4, No.~4 命题 5 的推论 3 以及同处, §1, No.~1 命题 4）；由此可见 \(N_{U}\) 是 \(\Gamma^{0}\) 的一个闭子集。此外，设 \(\varphi \in \mathcal{K}_{+}(G)\) 为一个函数，使得 \(\varphi(e) = 1\) 且在 G - U 上 \(\varphi(x) = 0\)；显然 \(\int \varphi(x)d\alpha(x) = 1\) 对所有 \(\alpha \in N_{U}\) 成立；于是 No.~1 的命题 2 表明 \(N_{U}\) 是一个紧集，这就证明了 a)。另一方面，设 V 为 \(e\) 在 G 中的一个相对紧开邻域，使得 \(\overline{V} \subset U\)，且设 \(\varphi \in \mathcal{K}_{+}(G)\)，其支集包含在 U 中且在 V 上 \(\varphi(x) = 1\)。则 \(\alpha(\varphi) = 1\)（对 \(\alpha \in N_{U}\)），因此存在 \(\alpha\) 在 N 中的一个邻域 W，使得 \(\beta(\varphi) < 2\)（对 \(\beta \in W\)）；于是显然 \(W \subset N_{V}\)，因此 \(N_{V}\) 是 \(N_{U}\) 的一个邻域。由于诸 \(N_{U}\) 覆盖 N，这就证明了 b)。最后，N 的每一紧子集 M 都包含在有限多个集合 \(N_{U_i}\)（\(1 \leqslant i \leqslant n\)）的并中，而由于 \(\bigcup_{i} N_{U_i} \subset N_{U}\)（其中 \(U = \bigcap_{i} U_i\)），这就证明了 c)。

推论. --- \(\Gamma^0\) 的子空间 N 是局部紧的。
% semantic-fragments: legacy-input-seam
\relax{}
